% Generated draft: requires mathematical and source review.
\documentclass[11pt]{article}
\usepackage[margin=25mm]{geometry}
\usepackage{fontspec}
\setmainfont{Latin Modern Roman}
\usepackage{amsmath,amssymb,mathtools,unicode-math}
\setmathfont{Latin Modern Math}
\usepackage{enumitem,needspace,xurl,hyperref}
\setlength{\parindent}{0pt}
\setlength{\parskip}{5pt}
\setlength{\emergencystretch}{3em}
\newcommand{\N}{\mathbb{N}}
\newcommand{\Z}{\mathbb{Z}}
\newcommand{\Q}{\mathbb{Q}}
\newcommand{\R}{\mathbb{R}}
\newcommand{\C}{\mathbb{C}}
\newcommand{\F}{\mathbb{F}}
\newcommand{\A}{\mathbb{A}}
\newcommand{\PP}{\mathbb{P}}
\newcommand{\E}{\mathbb{E}}
\newcommand{\eps}{\varepsilon}
\newcommand{\dd}{\,\mathrm{d}}
\begin{document}
\section*{UnsolvedMath ID: 30006341}
\textbf{Source problem code:} catalogue entry
\subsection{Definitions and mathematical statement}
For a finite group \(G\), define its minimal faithful permutation degree by
\[
\mu(G)=\min\{n\in\Z_{\geq 1}:\text{there exists an injective homomorphism }G\longrightarrow S_n\},
\]
where \(S_n\) is the symmetric group on \(n\) points. This definition allows actions with more than one orbit; the supplied passage does not explicitly specify whether transitivity is required.

The supplied source verification describes two candidate groups arising from conjugacy-class information for a prospective maximal subgroup of the Monster of shape
\[
2^{10+16}.O_{10}^{+}(2).
\]
Denote these candidates provisionally by \(G_1\) and \(G_2\). These symbols are placeholders: neither candidate is identified by a presentation, construction, or distinguishing extension data in the supplied record. The displayed shape alone does not specify either group or the relevant extension structure.

The reported conjecture is that, for each \(i\in\{1,2\}\), the value \(\mu(G_i)\) is approximately the quantity recorded as \(1.2\cdot107\). The numerical expression is preserved as supplied: its typography does not establish whether \(1.2\times10^7\) was intended. Moreover, ``about'' supplies no error tolerance, interval, or exact proposed value. Thus this is a contextual report of a conjecture, not a precise numerical equality or inequality. The passage also describes such permutation representations as difficult to handle computationally, without specifying a formal complexity claim.

The supplied record marks the statement as unrecoverable: the identities of the two candidates and a precise numerical target are missing. Consequently, a self-contained mathematical conjecture about these individual groups cannot be recovered from the supplied evidence.
\subsection{Short English statement}
The report conjectures an approximate minimal permutation degree for each of two unnamed candidate groups associated with a prospective Monster maximal subgroup. Their identities, the intended numerical notation, and the meaning of the approximation remain unresolved, so this draft preserves the context without asserting a precise conjecture.
\subsection{Sources}
\begin{itemize}
\item Local UnsolvedMath registry, permanent ID \(30006341\), title \emph{Unnamed Candidate Monster-Subgroup Extensions}, exact source problem code \texttt{catalogue entry}, in the Oberwolfach Reports Open Problems collection. The supplied record labels the problem open and its statement unrecoverable; these are supplied curation labels, not independently verified mathematical conclusions.
\item \emph{Computational Group Theory} (2025), Oberwolfach Reports, pp.\ 1389--1450, DOI \texttt{10.4171/owr/2025/27}. Supplied source URL and literature-source link labelled ``OWR report'': \url{https://doi.org/10.4171/owr/2025/27}. No individual proposer, precise passage location, or separate authoritative citation identifying the two groups is supplied.
\item Literature assessment dated \texttt{2026-08-22}: the supplied curation reports that the two candidate groups were constructed and their character tables computed, while their minimal-permutation-degree estimate remained conjectural. This is supplied dated curation, not an independently verified literature assessment. The background expressly requests independent verification of time-sensitive status and cited claims.
\item The supplied statement-verification note says that the official report describes two candidates for a prospective Monster maximal subgroup of shape \(2^{10+16}.O_{10}^{+}(2)\), but provides neither names nor presentations or distinguishing extension data. It records that no separate clean statement was promoted and that the original passage was retained for provenance. This account is attributed to the supplied verification note; the report has not been independently inspected here.
\end{itemize}
\end{document}
